\documentclass[a4paper]{article}

\title{Отчёт по лабораторной работе №1\\
<<Определение энтальпии нейтрализации сильной
кислоты сильным основанием>>}
\author{Можаров А.Р.}

\usepackage{cmap}
\usepackage[T2A]{fontenc}
\usepackage[utf8]{inputenc}
\usepackage[russian]{babel}
\usepackage{indentfirst}
\usepackage[hidelinks]{hyperref}
\usepackage{pgfplots}
\usepackage{amsmath}
\usepackage{amssymb}

\begin{document}

\maketitle

\subsection*{Цель работы}

Определение энтальпии нейтрализации сильной кислоты
сильным основанием по данным реакционной калориметрии

\subsection*{Оборудование}

\begin{enumerate}
    \item Модуль <<Термостат>>
          \begin{enumerate}
              \item Калориметр
              \item Центральный контроллер
              \item Термодатчик
              \item Нагревательный элемент
              \item Мешалка
              \item Насос
          \end{enumerate}
    \item Персональный компьютер
          \begin{enumerate}
              \item \texttt{elsms2.exe}
          \end{enumerate}
    \item Шприц
\end{enumerate}

\subsection*{Реагенты}

\begin{enumerate}
    \item 3 мл HCl
    \item 100 мл 0.15 $\dfrac{\text{моль}}{\text{л}}$ NaOH
    \item 2 г KCl
    \item 100 мл дистиллированной воды
\end{enumerate}

\section*{\centering Теоретическая часть}

Известно, что нейтрализация одного моля сильной одноосновной кислоты
сильным основанием в сильно разбавленных водных растворах
сопровождается одинаковым энергетическим эффектом,
равным
\begin{equation*}
    \Delta_\text{нейтр.}H = -13710\ \text{кал}\cdot\text{моль}^{-1}\ \text{при}\ T = 298.15\ K
\end{equation*}
Для экспериментального определения энтальпии нейтрализации проводят калориметрический опыт,
в котором в качестве калориметрической жидкости выступает раствор 0.15М раствор NaOH объёмом 100 мл.
Нейтрализацию проводят раствором выбранной кислоты заданной концентрации объёмом 3 мл,
который отбирают в специально подготовлены шприц.

Следует отметить, что при вливании раствора сильной кислоты в раствор сильного основания наряду
с процессом нейтрализации происходят ещё два процесса --- разбавление раствора кислоты основанием и раствора основания кислотой,
которые также сопровождаются энергетическими эффектами.
Если энтальпией второго процесса можно пренебречь в силу незначительности изменения объема раствора основания,
то энтальпия первого является весьма значительной величиной и в случае с HCl может составлять до 10\% от величины энтальпии нейтрализации.
В этой связи энтальпию разбавления HCl определяют в дополнительном калориметрическом опыте,
вливая 3 мл кислоты в 100 мл дистиллированной воды, выступающей в роли калориметрической жидкости.

Суммарную энтальпию в опыте и энтальпию разбавления раствора HCl рассчитывают по следующим уравнениям:
\begin{gather}
    -\Delta_\text{сумм.}H = (G \cdot c + W) \cdot \Delta T \\
    -\Delta_\text{разб.}H = (G \cdot c + W) \cdot \Delta T'
\end{gather}
где $G$ --- суммарная масса раствора в граммах,
$c$ --- удельная теплоёмкость раствора, которая принята величине, равной $0.98\ \text{кал}\cdot{г}^{-1}\cdot\text{К}^{-1}$,
$W$- энергетический эквивалент калориметрической системы,
$\Delta T$ и $\Delta T'$ --- изменения температуры в калориметрических опытах по определению энтальпией нейтрализации и разбавления, соответственно.

Вычитая энтальпию разбавления из суммарной энтальпии процесса, получают энтальпию нейтрализации, относящуюся ко взятому количеству кислоты:
\begin{equation}
    \Delta_\text{нейтр.}H = \Delta_\text{сумм.}H - \Delta_\text{разб.}H
\end{equation}
Значение мольной энтальпии нейтрализации кислоты получают делением величины $\Delta_\text{нейтр.}H$ на число молей кислоты,
равное $n = N \cdot V / 1000$ где $N$ и $V$ - нормальность и объёма (в мл) кислоты, соответственно.
Затем рассчитывают ошибку определения энтальпии нейтрализации кислоты.

Энергетический эквивалент калориметрической системы $W$ можно определить по следующему уравнению:
\begin{equation}
    -\Delta_\text{sol}H = \dfrac{M}{g} \cdot ((G + g) \cdot + W) \cdot \Delta T
\end{equation}
где  --- молярная энтальпия растворения KCl,
$G$ и $g$ --- массы воды и KCl в граммах,
$M$ --- молекулярная масса KCl,
$\Delta T$ --- изменение температуры в калориметрическом опыте.

\section*{\centering Практическая часть}

Для определения энтальпий процессов в данной работе используется
учебно-лабораторный комплекс «Химия», включающий модуль «Термостат».

Мы включили первый канал термодатчика и активировали контроллер.
Включили мешалку на необходимую интенсивность. Опыт проводился в три этапа.
В течении некоторого времени мы снимали показания температуры,
после в калориметрический стакан поместили навеску растворяемого вещества
(или заданный объём кислоты) после чего продолжили измерения.
После завершения калориметрического эксперимента,
сохранили данные и построили соответствующие графики.
Но для начала приведём табличные значения используемых величин:

\begin{table}[h]
    \begin{tabular}{|c|c|c|}
        \hline
        Обозначение          & Величина                          & Значение                                              \\
        \hline
        \hline
        $G$                  & Масса воды                        & 100 г                                                  \\
        \hline
        $g$                  & Масса KCl                         & 2 г                                                    \\
        \hline
        $c$                  & Удельная теплоёмкость раствора    & $0.98\ \text{кал}\cdot\text{г}^{-1}\cdot\text{К}^{-1}$ \\
        \hline
        $\Delta_\text{sol}H$ & Мольная энтальпия растворения KCl & 4197 кал/моль                                         \\
        \hline
        $M$                  & Молярная масса KCl                & 74.5 г/моль                                           \\
        \hline
    \end{tabular}
    \caption{Значения, взятые из справочника}
\end{table}

Добавили 2 г порошковая формы KCl в 100 мл дистиллированной воды
(растворение эталонного вещества с образованием раствора определённой концентрации).
Построили график и по его показаниям померили  (см. график 1)
Воспользовавшись табличными значениями величин нашли из формулы (4)
энергетический эквивалент калориметрической системы W.
\begin{equation}
    W = \dfrac{-\Delta_{sol}H \cdot g}{M \cdot \Delta T} - (G + g) \cdot c = 16.2\ \text{кал}\cdot\text{К}^{-1}
\end{equation}

\begin{figure}[h]
    \includegraphics[width=\textwidth]{1.jpg}
    \caption{Растворение КСl в дистиллированной воде}
\end{figure}

Добавили 3 мл кислоты HCl в 100 мл щёлочи NaOH (нейтрализация).
Аналогично прошлому опыту получили данные, которые преобразовали в график,
из которого определили  (см. график 2).
Воспользовавшись формулой (1) несложно найти:
\begin{equation}
    -\Delta_\text{сумм.}H = (G \cdot c + W) \cdot \Delta T_2 = 148.33\ \text{кал}
\end{equation}

\begin{figure}[h]
    \includegraphics[width=\textwidth]{2.jpg}
    \caption{Нейтрализация щёлочи NaOH кислотой HCl}
\end{figure}

Наконец рассмотрели добавление 3 мл кислоты HCl в 100 мл дистиллированной воды.
Стандартной процедурой определили  (см. график 3).
Энтальпия разбавления из формулы (2) определяется без особого труда:
\begin{equation}
    -\Delta_\text{разб.}H = (G \cdot c + W) \cdot \Delta T_3 = 19.40\ \text{кал}
\end{equation}

\begin{figure}[h]
    \includegraphics[width=\textwidth]{3.jpg}
    \caption{Разбавление кислоты HCl дистиллированной водой}
\end{figure}

Наконец получим искомую энтальпию нейтрализации сильной кислоты сильным основанием
из формулы (3), воспользовавшись результатами (6) и (7):
\begin{equation}
    \Delta_\text{нейтр.}H = \Delta_\text{сумм.}H - \Delta_\text{раств.}H = -167.73 \text{кал}
\end{equation}
Для получения мольной энтальпии нейтрализации нужно поделить $\Delta_\text{нейтр.}H$ на $n$
число молей кислоты:
\begin{equation*}
    \Delta_\text{нейтр.}H = -13997.5\ \text{кал}\cdot\text{моль}^{-1}
\end{equation*}
Эталонное значение взятое из справочника
\begin{equation*}
    \Delta_\text{нейтр.}H = -13351.48581\ \text{кал}\cdot\text{моль}^{-1}
\end{equation*}
Откуда отклонение от эталонного значения $4.6\%$.

\section*{\centering Вывод}

В ходе работы мы определили энтальпию нейтрализации сильной кислоты
сильным основанием по данным реакционной калориметрии.

Сравнили полученное значение с табличным и выявили,
что разница составила 4.6\% от эталонного значения.

\end{document}